\section{Banco de questões — Análise de Dados}

\subsection*{N1 — Fundamento competitivo}

\begin{questao}{\nivel{N1} 17.01 — Cebraspe/Certo ou Errado}
Em problema de classificação altamente desbalanceado, acurácia elevada pode coexistir com recall nulo para a classe de interesse.
\end{questao}

\begin{questao}{\nivel{N1} 17.02 — Cebraspe/Certo ou Errado}
Precisão responde à pergunta ``entre os casos previstos como positivos, quantos eram realmente positivos?'', enquanto recall mede a fração dos positivos reais detectados.
\end{questao}

\begin{questao}{\nivel{N1} 17.03 — Cebraspe/Certo ou Errado}
AUC-ROC elevada garante que as probabilidades previstas estejam bem calibradas.
\end{questao}

\begin{questao}{\nivel{N1} 17.04 — Cebraspe/Certo ou Errado}
Em validação de modelo temporal, divisão aleatória pode produzir avaliação excessivamente otimista se informações futuras vazarem para o treino.
\end{questao}

\begin{questao}{\nivel{N1} 17.05 — Cebraspe/Certo ou Errado}
Padronização deve ser ajustada apenas no conjunto de treino e depois aplicada aos conjuntos de validação e teste, para evitar leakage.
\end{questao}

\begin{questao}{\nivel{N1} 17.06 — Cebraspe/Certo ou Errado}
K-means é sensível à escala das variáveis e à inicialização dos centroides.
\end{questao}

\begin{questao}{\nivel{N1} 17.07 — Cebraspe/Certo ou Errado}
Lift maior que 1 em regra de associação indica associação positiva em relação à frequência-base, mas não demonstra causalidade.
\end{questao}

\begin{questao}{\nivel{N1} 17.08 — Cebraspe/Certo ou Errado}
Um outlier estatístico deve ser removido automaticamente antes de qualquer modelagem, pois representa necessariamente erro de dado.
\end{questao}

\begin{questao}{\nivel{N1} 17.09 — Cebraspe/Certo ou Errado}
RMSE penaliza erros grandes mais fortemente que MAE, em razão do quadrado dos resíduos antes da média e da raiz.
\end{questao}

\begin{questao}{\nivel{N1} 17.10 — Cebraspe/Certo ou Errado}
Correlação elevada entre duas variáveis é suficiente para concluir que uma causa a outra.
\end{questao}

\subsection*{N2 — Aplicação}

\begin{questao}{\nivel{N2} 17.11 — FCC/matriz de confusão}
Um classificador produziu 80 verdadeiros positivos, 20 falsos negativos e 40 falsos positivos. O recall da classe positiva é:
\begin{enumerate}[label=\Alph*)]
\item 50\%;
\item 66,7\%;
\item 80\%;
\item 90\%;
\item 100\%.
\end{enumerate}
\end{questao}

\begin{questao}{\nivel{N2} 17.12 — FGV/métrica}
Uma equipe de fiscalização só consegue analisar 100 alertas por semana. O custo principal é desperdiçar vagas com falsos positivos. A métrica a enfatizar entre os alertas selecionados é:
\begin{enumerate}[label=\Alph*)]
\item precisão;
\item recall apenas;
\item R\textsuperscript{2};
\item silhouette;
\item RMSE.
\end{enumerate}
\end{questao}

\begin{questao}{\nivel{N2} 17.13 — Cebraspe/Certo ou Errado}
Quando o custo de perder um caso positivo grave é elevado, recall pode ser mais relevante que precisão, dependendo da capacidade operacional e do contexto.
\end{questao}

\begin{questao}{\nivel{N2} 17.14 — FCC/F1}
Se precisão = 0,75 e recall = 0,60, o F1 é aproximadamente:
\begin{enumerate}[label=\Alph*)]
\item 0,45;
\item 0,60;
\item 0,67;
\item 0,75;
\item 0,90.
\end{enumerate}
\end{questao}

\begin{questao}{\nivel{N2} 17.15 — FGV/preprocessamento}
Uma equipe calcula média e desvio-padrão de todas as observações, inclusive do teste, antes de padronizar os dados. O problema é:
\begin{enumerate}[label=\Alph*)]
\item underflow;
\item data leakage;
\item normalização relacional;
\item deadlock;
\item somente overfitting de árvore.
\end{enumerate}
\end{questao}

\begin{questao}{\nivel{N2} 17.16 — Cebraspe/Certo ou Errado}
Árvores de decisão, em geral, são menos sensíveis ao escalonamento das variáveis do que algoritmos baseados em distância, como k-NN.
\end{questao}

\begin{questao}{\nivel{N2} 17.17 — FCC/categóricas}
Uma variável ``UF'' foi codificada como MA=1, PI=2, CE=3. Para algoritmo que interpreta distância numérica, o risco é:
\begin{enumerate}[label=\Alph*)]
\item criar ordem/distância artificial entre categorias nominais;
\item eliminar todas as categorias;
\item garantir melhor desempenho;
\item reduzir cardinalidade a zero;
\item tornar o modelo não supervisionado.
\end{enumerate}
\end{questao}

\begin{questao}{\nivel{N2} 17.18 — FGV/CRISP-DM}
Um projeto alcançou AUC 0,94, mas os auditores não conseguem usar o ranking porque o modelo foi treinado para um alvo sem relação com a decisão de fiscalização. A falha principal ocorreu em:
\begin{enumerate}[label=\Alph*)]
\item entendimento do negócio/definição do problema;
\item armazenamento;
\item visualização apenas;
\item normalização de banco;
\item clustering.
\end{enumerate}
\end{questao}

\begin{questao}{\nivel{N2} 17.19 — Cebraspe/Certo ou Errado}
No CRISP-DM, a passagem entre fases é iterativa; resultados da modelagem podem exigir retorno à preparação ou ao entendimento dos dados.
\end{questao}

\begin{questao}{\nivel{N2} 17.20 — FCC/ROC}
Ao reduzir o threshold de classificação, tende-se, em muitos modelos, a:
\begin{enumerate}[label=\Alph*)]
\item reduzir simultaneamente TPR e FPR sempre;
\item aumentar TPR e, em geral, também FPR;
\item manter matriz de confusão invariável;
\item eliminar falsos positivos;
\item tornar AUC igual a 1.
\end{enumerate}
\end{questao}

\begin{questao}{\nivel{N2} 17.21 — FGV/PR curve}
Em detecção de fraude com 0,2\% de positivos, a curva Precision-Recall pode ser particularmente informativa porque:
\begin{enumerate}[label=\Alph*)]
\item ignora a classe positiva;
\item enfatiza desempenho sobre a classe rara e os falsos positivos;
\item é equivalente a R\textsuperscript{2};
\item elimina a necessidade de threshold;
\item garante calibração.
\end{enumerate}
\end{questao}

\begin{questao}{\nivel{N2} 17.22 — Cebraspe/Certo ou Errado}
DBSCAN pode identificar ruído e clusters de forma não esférica, mas seu resultado depende de parâmetros de densidade e escala dos dados.
\end{questao}

\begin{questao}{\nivel{N2} 17.23 — FCC/regras de associação}
Se $P(B)=0,20$ e $P(B\mid A)=0,50$, o lift de $A\Rightarrow B$ é:
\begin{enumerate}[label=\Alph*)]
\item 0,10;
\item 0,40;
\item 1,0;
\item 2,5;
\item 5,0.
\end{enumerate}
\end{questao}

\begin{questao}{\nivel{N2} 17.24 — FGV/regressão}
Dois modelos têm MAE semelhante, mas o modelo X possui RMSE muito maior. Isso sugere que X:
\begin{enumerate}[label=\Alph*)]
\item possui alguns erros grandes que pesam fortemente no RMSE;
\item não possui outliers de erro;
\item é necessariamente melhor calibrado;
\item tem R\textsuperscript{2}=1;
\item não pode ser regressão.
\end{enumerate}
\end{questao}

\begin{questao}{\nivel{N2} 17.25 — Cebraspe/Certo ou Errado}
Em análise de texto, TF-IDF tende a aumentar o peso de termos importantes em um documento e relativamente raros no corpus, quando comparado à contagem bruta.
\end{questao}

\subsection*{N3 — Integração de concurso}

\begin{questao}{\nivel{N3} 17.26 — FGV/assertivas}
Considere:

I. Acurácia pode ser inadequada em classe rara.

II. AUC-ROC mede capacidade de ordenação ao longo de thresholds, não calibração diretamente.

III. Precision-Recall é inútil quando a classe positiva é rara.

Assinale a alternativa correta.
\begin{enumerate}[label=\Alph*)]
\item Apenas I.
\item Apenas II.
\item Apenas I e II.
\item Apenas II e III.
\item I, II e III.
\end{enumerate}
\end{questao}

\begin{questao}{\nivel{N3} 17.27 — Cebraspe/Certo ou Errado}
Um modelo com AUC elevada pode ter probabilidades sistematicamente superestimadas ou subestimadas; nesse caso, técnicas de calibração podem ser relevantes.
\end{questao}

\begin{questao}{\nivel{N3} 17.28 — FCC/leakage}
Um modelo de risco de contratos usa a variável ``quantidade final de aditivos'', calculada após encerramento do contrato, para decidir quais contratos em andamento serão fiscalizados. O problema é:
\begin{enumerate}[label=\Alph*)]
\item alta dimensionalidade;
\item leakage temporal;
\item falta de one-hot encoding;
\item clustering inadequado;
\item baixo recall necessariamente.
\end{enumerate}
\end{questao}

\begin{questao}{\nivel{N3} 17.29 — FGV/validação temporal}
Dados de 2019--2025 serão usados para prever casos de 2026. A estratégia de avaliação mais realista é:
\begin{enumerate}[label=\Alph*)]
\item embaralhar todos os anos e dividir aleatoriamente;
\item preservar ordem temporal, treinando no passado e validando/testando em períodos posteriores;
\item usar apenas 2025 para tudo;
\item calcular preprocessing em toda a série antes da divisão;
\item remover datas do dataset e dividir aleatoriamente.
\end{enumerate}
\end{questao}

\begin{questao}{\nivel{N3} 17.30 — Cebraspe/Certo ou Errado}
Se o processo de geração dos rótulos depende de quais casos foram historicamente auditados, o modelo pode reproduzir vieses de seleção existentes no processo de fiscalização.
\end{questao}

\begin{questao}{\nivel{N3} 17.31 — FCC/clustering}
K-means aplicado a duas variáveis, uma em reais de milhões e outra entre 0 e 1, sem escalonamento, tende a ser dominado por:
\begin{enumerate}[label=\Alph*)]
\item variável de maior escala;
\item variável binária apenas;
\item número de linhas;
\item target supervisionado;
\item matriz de confusão.
\end{enumerate}
\end{questao}

\begin{questao}{\nivel{N3} 17.32 — FGV/anomalias}
Um pagamento de R\$ 30 milhões é extremo globalmente, mas normal para grandes obras de infraestrutura. Isso ilustra:
\begin{enumerate}[label=\Alph*)]
\item necessidade de considerar anomalia contextual e porte do objeto;
\item que todo outlier é fraude;
\item que z-score sempre resolve;
\item que contexto deve ser removido;
\item que clustering é supervisionado.
\end{enumerate}
\end{questao}

\begin{questao}{\nivel{N3} 17.33 — Cebraspe/Certo ou Errado}
Uma regra de associação com alto confidence pode ser pouco informativa se o consequente já for muito frequente; o lift ajuda a comparar com essa frequência-base.
\end{questao}

\begin{questao}{\nivel{N3} 17.34 — FCC/visualização}
Um gráfico de barras começa o eixo Y em 98 em vez de 0 e compara valores 99 e 100, produzindo diferença visual enorme. O risco é:
\begin{enumerate}[label=\Alph*)]
\item distorção perceptiva da magnitude da diferença;
\item aumento de acurácia;
\item correção automática por legenda;
\item perda de dados do banco;
\item leakage.
\end{enumerate}
\end{questao}

\begin{questao}{\nivel{N3} 17.35 — FGV/SQL analítico}
Para ranquear fornecedores por valor dentro de cada órgão sem colapsar as linhas, a construção mais apropriada é:
\begin{enumerate}[label=\Alph*)]
\item função de janela com \texttt{PARTITION BY orgao};
\item apenas \texttt{GROUP BY fornecedor};
\item \texttt{DELETE};
\item \texttt{CROSS JOIN};
\item índice hash.
\end{enumerate}
\end{questao}

\begin{questao}{\nivel{N3} 17.36 — Cebraspe/Certo ou Errado}
Uma função de janela pode calcular agregados ou rankings sobre conjuntos relacionados sem reduzir o número de linhas como ocorre em um \texttt{GROUP BY} tradicional.
\end{questao}

\begin{questao}{\nivel{N3} 17.37 — FCC/Python}
Em pandas, um \texttt{groupby} seguido de agregação pode alterar a granularidade do dataset; por isso, ao fazer merge posterior, o analista deve verificar:
\begin{enumerate}[label=\Alph*)]
\item cardinalidade e chaves do merge;
\item apenas cor do gráfico;
\item memória da GPU;
\item versão do sistema operacional apenas;
\item somente nome das colunas.
\end{enumerate}
\end{questao}

\begin{questao}{\nivel{N3} 17.38 — FGV/planilhas}
Uma análise recorrente é realizada manualmente em planilha com fórmulas ocultas e cópias divergentes por e-mail. O principal risco é:
\begin{enumerate}[label=\Alph*)]
\item baixa reprodutibilidade, controle de versão e auditabilidade;
\item excesso de automação;
\item impossibilidade de calcular média;
\item falta de banco NoSQL;
\item uso de gráficos.
\end{enumerate}
\end{questao}

\begin{questao}{\nivel{N3} 17.39 — Cebraspe/Certo ou Errado}
Um embedding aproxima objetos em espaço vetorial por relações aprendidas, mas similaridade vetorial não equivale automaticamente a identidade ou causalidade.
\end{questao}

\begin{questao}{\nivel{N3} 17.40 — FCC/avaliação de anomalia}
Um detector sinaliza 500 contratos como anômalos. Após revisão humana, apenas 30 são relevantes. Antes de concluir que o algoritmo é inútil, deve-se avaliar:
\begin{enumerate}[label=\Alph*)]
\item threshold, objetivo, capacidade de investigação e custo de falsos positivos;
\item somente CPU;
\item número de colunas;
\item formato CSV;
\item apenas o nome do algoritmo.
\end{enumerate}
\end{questao}

\subsection*{N4 — Auditoria / avançado}

\begin{questao}{\nivel{N4} 17.41 — Cebraspe/caso integrado}
Um modelo para selecionar licitações tem AUC 0,93, foi treinado com divisão aleatória de 2018--2025 e usa atributos disponíveis apenas após a assinatura do contrato. Julgue o item:

O AUC elevado é evidência suficiente de que o modelo pode ser usado na fase pré-contratual.
\end{questao}

\begin{questao}{\nivel{N4} 17.42 — FGV/auditoria analítica}
Um dashboard oficial mostra total 8\% acima da fonte. O pipeline une contratos e pagamentos em relação 1:N. O primeiro procedimento é:
\begin{enumerate}[label=\Alph*)]
\item aplicar \texttt{DISTINCT} indiscriminadamente;
\item mapear granularidade, cardinalidade dos joins e reconciliar totais por etapa;
\item excluir registros grandes;
\item trocar ferramenta de BI;
\item aumentar memória.
\end{enumerate}
\end{questao}

\begin{questao}{\nivel{N4} 17.43 — FCC/viés de seleção}
Somente contratos já auditados possuem rótulo confiável de irregularidade. Um modelo treinado nesses dados pode:
\begin{enumerate}[label=\Alph*)]
\item reproduzir o padrão histórico de seleção e ignorar áreas nunca fiscalizadas;
\item eliminar qualquer viés por ser supervisionado;
\item garantir fairness;
\item dispensar validação externa;
\item produzir causalidade.
\end{enumerate}
\end{questao}

\begin{questao}{\nivel{N4} 17.44 — Cebraspe/Certo ou Errado}
Em auditoria de modelo, é insuficiente verificar apenas métrica global; desempenho por subgrupo, drift, disponibilidade das features no momento da decisão e processo de supervisão humana podem ser materialmente relevantes.
\end{questao}

\begin{questao}{\nivel{N4} 17.45 — FGV/reprodutibilidade}
Um relatório de fiscalização foi produzido por notebook sem versão, com dados sobrescritos e sem registro de parâmetros. A principal recomendação é:
\begin{enumerate}[label=\Alph*)]
\item salvar somente o PDF final;
\item versionar código, parâmetros, referências dos dados e ambiente de execução;
\item migrar obrigatoriamente para SAS;
\item remover comentários;
\item impedir qualquer notebook.
\end{enumerate}
\end{questao}

\begin{questao}{\nivel{N4} 17.46 — Cebraspe/Certo ou Errado}
Um modelo pode manter métricas técnicas estáveis e, ainda assim, perder valor operacional se a capacidade da equipe para investigar alertas mudar significativamente.
\end{questao}

\begin{questao}{\nivel{N4} 17.47 — FCC/calibração}
Um modelo ordena bem os riscos, mas casos com probabilidade prevista de 80\% apresentam frequência real de 45\%. O problema é principalmente de:
\begin{enumerate}[label=\Alph*)]
\item calibração;
\item clustering;
\item normalização de banco;
\item cardinalidade de SQL;
\item encoding de texto apenas.
\end{enumerate}
\end{questao}

\begin{questao}{\nivel{N4} 17.48 — FGV/governança}
Uma área altera o threshold do modelo para aumentar o número de alertas, mas não registra a mudança. Qual risco surge?
\begin{enumerate}[label=\Alph*)]
\item perda de rastreabilidade, comparabilidade e controle da decisão analítica;
\item aumento automático de recall sem custo;
\item melhoria garantida de precisão;
\item nenhuma mudança, pois o modelo é o mesmo;
\item somente aumento de storage.
\end{enumerate}
\end{questao}

\begin{questao}{\nivel{N4} 17.49 — Cebraspe/Certo ou Errado}
Explicabilidade local de uma predição não substitui validação global do modelo nem análise de qualidade dos dados usados para produzi-la.
\end{questao}

\begin{questao}{\nivel{N4} 17.50 — FGV/matriz de achado}
Um painel de risco não registra versão do modelo, data do dataset, threshold ou responsáveis pela validação. Qual achado é mais adequado?
\begin{enumerate}[label=\Alph*)]
\item ``Machine learning não deve ser usado no setor público'';
\item ``A ausência de versionamento e governança dos parâmetros reduz reprodutibilidade e accountability das decisões; recomenda-se registry, documentação, aprovação e trilha de execução'';
\item ``Basta aumentar o AUC'';
\item ``O painel deve ser substituído por planilha'';
\item ``Não há problema se os resultados parecem plausíveis''.
\end{enumerate}
\end{questao}
